% ============================================================
%  Setmana 4 — Unitat 7 · Límits i continuïtat
%  Exercicis: 112, 113, 114, 43, 120
% ============================================================

\setmana{4}{Unitat 7 \textperiodcentered{} Límits i continuïtat}{11 d'octubre de 2026}%
        {112, 113, 114, 43 i 120}

\begin{errada}{Exercici 43}
  \ondiu{$g(2)-4<0$}
  \hadedir{$g(2)=-4<0$}
  \motiu{Hi falta el signe igual: $g(2)=3^{-2+2}-5=1-5=-4$.
    El teorema de Bolzano necessita que $g$ tingui signes diferents als
    extrems de l'interval, i els té: $g(0)=4>0$ i $g(2)=-4<0$.}
\end{errada}

%\senseerrades{112, 113, 114 i 120}

% ---------- Observacions (només surten amb \obstrue) ----------
\ifobs
\bigskip
\noindent
Dues observacions que no són errades, però que poden confondre:

\begin{observacio}{Exercici 113}
  El valor $x=-3$ és una arrel exacta, perquè
  $x^3+6x^2+12x+9=(x+2)^3+1$. Si proves amb $x=-3$ trobaràs $f(-3)=0$ en lloc
  d'un canvi de signe: no t'has equivocat, has trobat l'arrel. Qualsevol
  interval de longitud $2$ que la contingui, com ara $(-4,-2)$, respon la
  pregunta.
\end{observacio}

\begin{observacio}{Exercici 120}
  L'enunciat anomena $f(x)=x+1$, i la solució torna a fer servir la lletra $f$
  per a una funció diferent: la resta de les dues funcions que es tallen (per
  exemple, $f(x)=x+1-e^{x-1}$ a l'apartat a). El punt on aquesta resta val
  zero és l'abscissa del punt de tall. D'altra banda, l'enunciat només demana
  demostrar que hi ha un punt de tall; la solució va més enllà i l'aproxima
  amb un error menor que una dècima.
\end{observacio}
\fi
